\subsubsection{Equivalent Reflection Point (ERP)}
\label{subsubsec:equivalent-reflection-points}

In complex environments, an electromagnetic wave often undergoes multiple interactions such as specular reflections, edge diffractions, or diffuse scattering before returning to the receiver.
Tracing every interaction sequence requires prior knowledge of the object's exact geometry and material properties.
Instead, the core idea behind detection in this framework is to abstract these potentially complex multipath trajectories into an Equivalent Reflection Point (ERP).
The ERP combines the total delay and direction change of the entire signal path into a single 3D point, removing the need to find every individual bounce along the way.

In our proposed framework, a monostatic ISAC setup is used where the transmitter and receivers are collocated.
Each remaining path $\ell$ is mapped to a single ERP in 3D space:
\begin{equation}
    \mathbf{p}_\ell^{\mathrm{ERP}}
    = \mathbf{p}_{\mathrm{gNB}}
    + \!\left(\frac{L_\ell}{2} + \delta\right)\mathbf{d}_\ell,
    \label{eq:erp}
\end{equation},
where $\mathbf{p}_{\mathrm{gNB}}$ is the gNB spatial location, and $\mathbf{d}_\ell$ is the unit direction vector of the path.
To account for the round-trip distance in monostatic case, we divide the path delay by 2($L_\ell/2$).
$\delta$, a small calibration offset is added to mitigates the bias introduced by reflection on an object's exterior surface rather than precisely at its geometric centre.
By transforming raw angle-delay measurements into 3D coordinates, the ERP allows to reconstruct the target's shape and locate its position.